%% Literature Survey — all references preserved.

\section{Literature Survey}

\subsection{Classical H(div) Spaces and Mimetic Methods}

The classical $H(\mathrm{div})$-conforming finite elements are the
\emph{Raviart--Thomas} (RT) family \cite{RaviartThomas1977} and the
\emph{Brezzi--Douglas--Marini} (BDM) family \cite{bdm1985}.  On a triangle
$T$, the RT$_k$ space is $\mathrm{RT}_k(T) = [\mathbb{P}_k(T)]^2 +
\mathbb{P}_k(T)\,\mathbf{x}$, with DOFs consisting of edge normal-trace
moments up to degree $k$ and interior moments of the vector field
\cite{BrezziFortin1991,BoffiBrezziFortin2013}.  N\'{e}d\'{e}lec's face
elements \cite{Nedelec1980,nedelec1986} generalize these to three dimensions.
On quadrilaterals, the Piola transform $\hat{\mathbf{v}} \mapsto
(J\hat{\mathbf{v}})/|\det J|$ maps $H(\mathrm{div})$ functions from a
reference square to the physical element; Boffi and Gastaldi
\cite{boffi2002remarks} showed that maintaining approximation order on
distorted quadrilaterals requires the full $\mathbb{Q}_k$ space in the
Piola-transformed formulation.

For general polygonal and polyhedral cells, \emph{mimetic finite difference}
(MFD) methods construct local inner products that mimic the continuous ones.
The foundational work of Brezzi, Lipnikov, Shashkov, and Simoncini
\cite{brezzi2005family} introduced a family of MFD methods for arbitrary
polygonal meshes, decomposing the local stiffness matrix into a consistency
part (exact for linear polynomials) and a positive semidefinite stability
correction.  The comprehensive survey \cite{lipnikov2014mimetic} showed how
this philosophy -- encoding both primal and dual DOFs -- produces discrete
operators satisfying exact sequences analogous to the de Rham complex,
with the discrete divergence theorem holding to machine precision.
Higher-order MFD on polygons was pursued in \cite{beiraoMFDhigh2014},
achieving $O(h^{k+1})$ convergence in $L^2$ using projector-based
constructions.  For curvilinear grids, \cite{lipnikov2014curvilinear} showed
that the method remains optimally convergent when the geometry is
approximated to a sufficiently high order -- a theme central to spherical
manifold extensions.

\subsection{Virtual Element Methods and H(div) VEM}

The \emph{Virtual Element Method} (VEM) of Beir\~{a}o da Veiga et al.\
\cite{veiga2013basic,BeiraoDaVeiga2013} extends MFD within a Galerkin
framework by using \emph{virtual} basis functions -- implicitly defined as
local PDE solutions -- that are never explicitly computed.  The
$H(\mathrm{div})$-conforming VEM \cite{veigaHdivVEM2016,BeiraoDaVeiga2016,%
BeiraoDaVeigaMixed2016} constructs spaces $\mathbf{V}_h \subset
H(\mathrm{div};\Omega)$ on polygonal meshes with DOFs consisting of
edge normal-trace moments (up to degree $k$), interior moments of the
vector field, and interior moments of the divergence.  The resulting mixed
pair satisfies a discrete inf-sup condition uniformly in $h$ on polygonal
meshes \cite{BrezziFalkMarini2014}.  The three-dimensional polyhedral
extension \cite{veiga3DVEM2017} proves exactness of the VEM complex,
required to avoid spurious pressure modes.  Serendipity VEM
\cite{brezziSerendipityVEM2017} reduces internal DOFs using exact-sequence
arguments, significantly lowering the cost of high-order schemes.

\subsection{Hybrid High-Order and Discrete de Rham Methods}

The \emph{Hybrid High-Order} (HHO) method of Di Pietro and Ern
\cite{dipietroHHO2015} achieves arbitrary-order accuracy on polytopal meshes
by attaching DOFs to both cell interiors and faces, with a local
reconstruction operator lifting face DOFs to a higher-degree polynomial.
For $H(\mathrm{div})$ reconstruction, \cite{dipietroHHOdiv2017} showed that
HHO's local reconstruction can be cast as an $H(\mathrm{div})$ projection,
closely related to HDG \cite{cockburnHDG2009} and the weak Galerkin method
\cite{dipietroBridge2016}.

The \emph{Discrete de Rham} (DDR) complex of Di Pietro, Droniou, and
Rapetti \cite{dipietroDDR2021} represents the current state of the art.
For arbitrary order $k\ge0$ on polyhedral meshes, the DDR complex constructs
a commuting diagram:
\[
\begin{tikzcd}
  \mathcal{P}^{k+1} \arrow[r,"\nabla"] \arrow[d] &
  \mathcal{P}^k \arrow[r,"\nabla\times"] \arrow[d] &
  \mathcal{P}^k \arrow[r,"\nabla\cdot"] \arrow[d] &
  \mathcal{P}^{k-1} \arrow[d] \\
  \underline{X}^{k+1}_h \arrow[r,"\underline{G}_h"] &
  \underline{X}^{k,\mathrm{curl}}_h \arrow[r,"\underline{C}_h"] &
  \underline{X}^{k,\mathrm{div}}_h \arrow[r,"\underline{D}_h"] &
  X^{k-1}_h
\end{tikzcd}
\]
with vertex, edge, face, and cell DOFs linked by commuting interpolation
operators.  Optimal convergence ($O(h^{k+1})$ in $L^2$, $O(h^k)$ in the
energy norm) is proved via a local potential reconstruction technique.

\subsection{H(div) Schemes on Spherical Manifolds}

Constructing $H(\mathrm{div})$-conforming schemes on the sphere introduces
additional challenges: geometric nonlinearity from curved edges, the tangential
constraint $\mathbf{u}\cdot\hat{r}=0$, nontrivial topology, and panel
interfaces on cubed-sphere and icosahedral grids.  Thuburn and Cotter
\cite{thuburn2012primal} developed a primal-dual mimetic scheme for the
rotating shallow water equations on Voronoi/Delaunay spherical meshes,
storing only normal velocity components on primal edges with the divergence
computed exactly via the discrete divergence theorem.  The MPAS atmospheric
model \cite{ringler2010unified} demonstrated that such Voronoi-based mimetic
schemes produce stable, long-term integrations with excellent energy
conservation.

The curved MFD approach of \cite{lipnikovCurvedMFD2023} extended MFD to
meshes with curved faces by isoparametric mapping, showing that geometric
approximation order $p\ge k+1$ is required for $O(h^k)$ convergence.
For the VEM setting, \cite{veigaCurvedVEM2019} introduced VEM on curved
triangular surfaces, proving that the virtual functions automatically adapt
to manifold geometry.  The DDR complex on manifolds
\cite{dipietroManifold2024} achieves arbitrary-order convergence on curved
polyhedral surface meshes by using tangential polynomial DOFs and the
manifold Hodge--Laplacian as the local operator.

\subsection{Harmonic Bases and Vector Spherical Harmonics}

\emph{Harmonic polynomials} -- solutions of the Laplace equation -- offer a
natural basis for $H(\mathrm{div})$ reconstruction.  Since
$\nabla\cdot(\nabla p) = \Delta p = 0$ for harmonic $p$, any gradient of a
harmonic polynomial is automatically divergence-free, and the
Helmholtz--Hodge decomposition $\mathbf{u} = \nabla\phi +
\nabla\times\psi + \mathbf{u}_0$ can be approximated by harmonic
polynomials for $\phi$ and $\psi$.  In 2D, the real and imaginary parts of
complex polynomials $(x+iy)^k$ provide a complete harmonic family.
Cockburn, Gopalakrishnan, and Guzm\'{a}n \cite{cockburnHarmonic2010} used
harmonic polynomial spaces for projection-based $H(\mathrm{div})$ and
$H(\mathrm{curl})$ interpolation, achieving commuting projections with
optimal estimates independent of element geometry.

On the sphere, the natural analogues are \emph{vector spherical harmonics}
(VSH): $\mathbf{V}_\ell^m = \nabla_s Y_\ell^m$ (irrotational) and
$\mathbf{W}_\ell^m = \hat{r}\times\nabla_s Y_\ell^m$ (solenoidal), forming
a complete orthonormal basis for tangential vector fields
\cite{wangVSH2017}.  For $hp$-adaptivity, hierarchical $H(\mathrm{div})$
bases \cite{zaglmayrHierarchical2006} use recursive edge/face/interior
splitting with well-conditioned stiffness matrices.  The LDG approach of
Cockburn and Shu \cite{cockburnLDG1998} provides $O(h^{k+1})$ convergence
for convection-dominated problems using $H(\mathrm{div})$-conforming
numerical fluxes, with the discrete Poincar\'{e} constant remaining bounded
as $k\to\infty$ \cite{ciarletStability2002}.

\subsection{Generalized Barycentric Coordinates and Whitney Forms}

Wachspress coordinates \cite{Wachspress1975} are rational generalized
barycentric coordinates on convex polygons, part of a wider family including
mean-value, harmonic, Sibson, and maximum-entropy coordinates
\cite{FloaterHormannKos2006,WarrenSchaeferHiraniDesbrun2007,Sibson1981,%
HormannSukumar2008}.  Sukumar and Tabarraei \cite{SukumarTabarraei2004}
demonstrated their use for conforming $H^1$ polygonal elements, and
Gillette, Rand, and Bajaj \cite{GilletteRandBajaj2012,RandGilletteBajaj2013}
established interpolation estimates and conditioning properties.  These works
support scalar $H^1$ interpolation but do not produce edge-flux-conservative
$H(\mathrm{div})$ spaces.  The bridge is provided by generalized Whitney-form
constructions \cite{GilletteRandBajaj2016}, which build $H(\mathrm{curl})$,
$H(\mathrm{div})$, and scalar families from barycentric coordinates on convex
polytopes, recovering classical RT, BDM, and N\'{e}d\'{e}lec behavior on
simplices.

\subsection{Conservative Remapping}

Conservative remapping between nonmatching grids has a separate mature
literature.  Integral remapping for ALE methods goes back to Dukowicz and
Kodis \cite{DukowiczKodis1987}.  Jones \cite{Jones1999} described first- and
second-order conservative remapping weights for spherical climate grids.
Ullrich and Taylor \cite{UllrichTaylor2015} and Ullrich, Devendran, and
Johansen \cite{UllrichDevendranJohansen2016} developed arbitrary-order
conservative and consistent linear maps, implemented in TempestRemap.  Those
methods transfer scalar cell averages; the present method transfers directed
edge fluxes, which is an $H(\mathrm{div})$-type transfer problem.

\subsection{Connection to the Present Work}

The Perot--Chartrand construction used here \cite{PerotChartrand2021}
belongs to the compatible polygonal family: it reconstructs a local vector
field from edge fluxes using a constant-divergence term and harmonic
gradients.  Our higher-order extension (Section~3) enriches this with a
split-moment basis decomposition into polynomial-divergence, harmonic-gradient,
and divergence-free completion modes, using the DOF structure
of $H(\mathrm{div})$ VEM \cite{veigaHdivVEM2016} (edge normal-trace moments
plus optional cell vector moments) with the harmonic basis stability of
\cite{cockburnHarmonic2010}.  The spherical extension uses gnomonic charts with
a contravariant Piola mapping, consistent with the curved-element analysis of
\cite{lipnikovCurvedMFD2023}, and degree elevation to capture the leading
rational correction from the Piola map.  The geometric tools are standard:
Voronoi diagrams \cite{Aurenhammer1991} and Sutherland--Hodgman polygon
clipping \cite{SutherlandHodgman1974}.
